\section{Homework Week 11}

\begin{exercise}
    \begin{enumerate}
        \item[(i)] Let \((X,\rho)\) be a transitive \(\operatorname{Set}\)-representation of \(G\).
        For \(x\in X\), let \(G_x\) denote the fixator of \(x\) in \(G\). Then the map
        \[
            G/G_x \longrightarrow X,
            \qquad
            gG_x \longmapsto \rho(g)(x)
        \]
        defines an isomorphism
        \[
            (G/G_x,\rho_{G_x}) \xrightarrow{\sim} (X,\rho).
        \]

        \item[(ii)] For subgroups \(H\) and \(H'\) of \(G\), the representations
        \[
            (G/H,\rho_H)
            \qquad\text{and}\qquad
            (G/H',\rho_{H'})
        \]
        are isomorphic if and only if \(H\) and \(H'\) are conjugate in \(G\).
    \end{enumerate}
\end{exercise}

\begin{exercise}
    Assume that \(G\) is the cyclic additive group \(\mathbb{Z}/p\mathbb{Z}\), where
    \(p\) is a prime number, and let \(k\) be a field of characteristic \(p\).

    Consider the morphism
    \[
        G \longrightarrow \operatorname{GL}_2(k),
        \qquad
        \lambda \longmapsto
        \begin{pmatrix}
            1 & \lambda \\
            0 & 1
        \end{pmatrix}.
    \]
    Prove that the line generated by the first standard basis vector of \(k^2\) is
    \(G\)-stable but has no \(G\)-stable complement.
\end{exercise}

\begin{exercise}
    Let \((M,\rho)\) and \((N,\sigma)\) be two \(k\)-representations of \(G\).
    Fix bases \((e_i)_{1\leq i\leq m}\) and \((f_j)_{1\leq j\leq n}\) of \(M\) and \(N\),
    respectively. Then
    \[
        (e_i\otimes f_j)_{1\leq i\leq m,\ 1\leq j\leq n}
    \]
    is a basis of \(M\otimes_k N\). Denote by
    \[
        (m,\rho),\qquad (n,\sigma),\qquad (mn,\rho\otimes\sigma)
    \]
    the matrix representations associated with
    \[
        (M,\rho),\qquad (N,\sigma),\qquad (M\otimes_k N,\rho\otimes\sigma),
    \]
    respectively.

    Determine the relation between \((\rho\otimes\sigma)(g)\) and
    \(\rho(g),\sigma(g)\).
\end{exercise}

\begin{exercise}
    Let \((M,\rho)\) be a \(k\)-representation of \(G\). Fix a basis
    \((e_i)_{1\leq i\leq m}\) of \(M\). Let \((M^*,\rho^*)\) be the contragredient
    representation of \((M,\rho)\). Denote by
    \[
        (m,\rho),\qquad (m,\rho^*)
    \]
    the matrix representations associated with
    \[
        (M,\rho),\qquad (M^*,\rho^*).
    \]

    Determine the relation between \(\rho(g)\) and \(\rho^*(g)\).
\end{exercise}

\begin{exercise}
    \begin{enumerate}
        \item[(i)] Let \(G\) be a finite group acting on a finite set \(\Omega\). Let
        \(\operatorname{Fix}^G(\Omega)\) denote the set of fixed points of \(\Omega\)
        under \(G\), and let \(\Omega^\#\) denote the set of orbits of \(G\) on
        \[
            \Omega\setminus \operatorname{Fix}^G(\Omega).
        \]
        Prove that
        \[
            |\Omega|
            =
            |\operatorname{Fix}^G(\Omega)|
            +
            \sum_{\omega\in \Omega^\#} |\omega|.
        \]

        From now on, we assume that \(p\) is a prime number and that \(G\) is a
        nontrivial \(p\)-group.

        \item[(ii)] Assume now that \(|\Omega|\) is a power of \(p\) different from \(1\).
        Prove that \(|\operatorname{Fix}^G(\Omega)|\) is divisible by \(p\).

        \item[(iii)] Let \(k\) be a commutative field of characteristic \(p\). Let \(M\)
        be a \(kG\)-module. Prove that
        \[
            \operatorname{Fix}^G(M) \neq 0.
        \]
        \textup{(HINT.} 
        \begin{enumerate}
            \item[(a)] Let \((e_i)_{i=1,\ldots,r}\) be a \(k\)-basis of \(M\).
            Prove that the abelian group generated by
            \[
                (ge_i)_{i=1,\ldots,r,\ g\in G}
            \]
            is an \(\mathbb{F}_pG\)-module.
            \item[(b)] Apply question \textup{(2)} above to conclude.
        \end{enumerate}
        \textup{)}

        \item[(iv)] Prove that, up to isomorphism, the trivial representation of \(G\)
        is the unique irreducible \(kG\)-module.
    \end{enumerate}
\end{exercise}